\question{}

\begin{ابجد}
\item طبق الگوریتم \lr{Value Iteration} در هر مرحله ارزش حالت‌ها را طبق رابطه‌ی زیر به‌روز می‌کنیم:
\begin{align*}
    V_{new}(s) = \max_a \sum_{s', r'} p(s', r \mid s,a) [r + \gamma V(s')]
\end{align*}
بنابراین مرحله‌ی اول این الگوریتم برای حالت‌ها به این صورت می‌باشد (خانه‌ها را با $(row, column)$ نمایش داده و قرارداد می‌کنیم که $V_i(s)$ ارزش خانه $s$ در سیاست $\pi_i$ هست):
\begin{align*}
    V_{1}(1, 1) & = \max \{ (0.05(0 + 0.9(0)) + 0.05(0 + 0.9(0)) + 0.9(0 + 0.9(0))),                    \\
                & \phantom{{}= \max \{ } (0.05(0 + 0.9(0)) + 0.05(0 + 0.9(0)) + 0.9(0 + 0.9(0))),       \\
                & \phantom{{}= \max \{ } (0.05(0 + 0.9(0)) + 0.05(0 + 0.9(0)) + 0.9(0 + 0.9(0))),       \\
                & \phantom{{}= \max \{ } (0.05(0 + 0.9(0)) + 0.05(0 + 0.9(0)) + 0.9(0 + 0.9(0))) \} = 0 \\[1em]
    %
    V_{1}(1, 2) & = \max \{ (0.05(-10 + 0.9(0))) + 0, (0.05(-10 + 0.9(0))) + 0,                         \\
                & \phantom{{}= \max \{ } (0.9(-10 + 0.9(0))) + 0, 0 \} = 0                              \\[1em]
    %
    V_{1}(1, 3) & = 0                                                                                   \\[1em]
    %
    V_{1}(2, 1) & = \max \{ 0, 0, 0, 0 \} = 0                                                           \\[1em]
    %
    V_{1}(2, 2) & = \max \{ (0.05(+10 + 0.9(0))) + 0, (0.05(+10 + 0.9(0))) + 0,                         \\
                & \phantom{{}= \max \{ } (0.9(+10 + 0.9(0))) + 0, 0 \} = 9                              \\[1em]
    %
    V_{1}(2, 3) & = 0
\end{align*}

و در \lr{iteration} دوم داریم:
\begin{align*}
    V_{2}(1, 1) & = \max \{ 0, 0, 0, 0 \} = 0                                                & \\[1em]
    %
    V_{2}(1, 2) & = \max \{ 0.05(-10 + 0.9(0)) + 0, 0.05(0 + 0.9(9)) + 0,                    & \\
                & \phantom{{}= \max \{ } 0.05(-10 + 0.9(0)) + 0.9(0 + 0.9(9)) + 0,           & \\
                & \phantom{{}= \max \{ } 0.05(0 + 0.9(9)) + 0.9(-10 + 0.9(0)) + 0 \} = 6.79  & \\[1em]
    %
    V_{2}(1, 3) & = 0                                                                        & \\[1em]
    %
    V_{2}(2, 1) & = \max \{ 0.05(0 + 0.9(9)) + 0, 0.05(0 + 0.9(9)) + 0,                      & \\
                & \phantom{{}= \max \{ } 0.9(0 + 0.9(9)) + 0, 0 \} = 7.29                    & \\[1em]
    %
    V_{2}(2, 2) & = \max \{ 0.05(+10 + 0.9(0)) + 0, 0.05(0 + 0.9(9)) + 0,                    & \\
                & \phantom{{}= \max \{ } 0.05(+10 + 0.9(0)) + 0.9(0 + 0.9(9)) + 0,           & \\
                & \phantom{{}= \max \{ } 0.05(0 + 0.9(9)) + 0.9(+10 + 0.9(0)) + 0 \} = 9.405 & \\[1em]
    %
    V_{2}(2, 3) & = 0                                                                        &
\end{align*}


\item در الگوریتم \lr{Temporal Difference Learning}، به‌روزرسانی‌ها طبق این رابطه انجام می‌شود:
\begin{align*}
    V_{new}^{\pi}(s) = (1 - \alpha)V^{\pi}(s) + (\alpha)(R(s, \pi(s), s') + \gamma V^{\pi}(s'))
\end{align*}
بنابراین در تجربه‌ی $(1,1) \rightarrow (1,2) \rightarrow (1,3)$ داریم:
\begin{align*}
    (1,1) \rightarrow (1,2): & \quad V_{new}^{\pi}(1,1) = (0.9)(0) + (0.1)(0 + 0.9(0)) = 0    \\
    (1,2) \rightarrow (1,3): & \quad V_{new}^{\pi}(1,2) = (0.9)(0) + (0.1)(-10 + 0.9(0)) = -1
\end{align*}
توجه داریم که برای تجربه‌ی دوم، باید از این مقادیر استفاده کنیم. در این تجربه داریم:
\begin{align*}
    (1,1) \rightarrow (1,2): & \quad V_{new}^{\pi}(1,1) = (0.9)(0) + (0.1)(0 + 0.9(-1)) = -0.09 \\
    (1,2) \rightarrow (2,2): & \quad V_{new}^{\pi}(1,2) = (0.9)(-1) + (0.1)(0 + 0.9(0)) = -0.9  \\
    (2,2) \rightarrow (2,3): & \quad V_{new}^{\pi}(2,2) = (0.9)(0) + (0.1)(+10 + 0.9(0)) = +1
\end{align*}

\item به‌روزرسانی در الگوریتم \lr{SARSA} به این صورت می‌باشد:
\begin{align*}
    Q_{new}(s_t, a_t) = Q(s_t, a_t) + \alpha [R_{t+1} + \gamma Q(s_{t+1}, a_{t+1}) - Q(s_t, a_t)]
\end{align*}
مسیر اول:
\begin{align*}
    (1,1) \rightarrow (1,2): & \quad Q_1((1,1), \text{\lr{right}}) = 0 + (0.1) [0 + 0.9(0) - 0] = 0    \\
    (1,2) \rightarrow (1,3): & \quad Q_1((1,2), \text{\lr{right}}) = 0 + (0.1) [-10 + 0.9(0) - 0] = -1
\end{align*}
و مسیر دوم:
\begin{align*}
    (1,1) \rightarrow (1,2): & \quad Q_1((1,1), \text{\lr{right}}) = 0 + (0.1) [0 + 0.9(-1) - 0] = -0.09   \\
    (1,2) \rightarrow (2,2): & \quad Q_1((1,2), \text{\lr{right}}) = -1 + (0.1) [0 + 0.9(0) - (-1)] = -0.9 \\
    (2,2) \rightarrow (2,3): & \quad Q_1((2,2), \text{\lr{right}}) = 0 + (0.1) [+10 + 0.9(0) - 0] = +1
\end{align*}

\item به‌روز‌رسانی در الگوریتم \lr{Q-Learning} به این صورت می‌باشد:
\begin{align*}
    Q_{new}(s_t, a_t) = Q(s_t, a_t) + \alpha [R_{t+1} + \gamma \max_a Q(s_{t+1}, a) - Q(s_t, a_t)]
\end{align*}
بنابراین این الگوریتم برای مسیر اول به این شکل خواهد بود:
\begin{align*}
    (1,1) \rightarrow (1,2): & \quad Q_1((1,1), \text{\lr{right}}) = 0 + (0.1) [0 + 0.9\max\{0,0,0,0\} - 0] = 0    \\
    (1,2) \rightarrow (1,3): & \quad Q_1((1,2), \text{\lr{right}}) = 0 + (0.1) [-10 + 0.9\max\{0,0,0,0\} - 0] = -1
\end{align*}
و برای مسیر دوم:
\begin{align*}
    (1,1) \rightarrow (1,2): & \quad Q_1((1,1), \text{\lr{right}}) = 0 + (0.1) [0 + 0.9\max\{0,-1,0,0\} - 0] = 0       \\
    (1,2) \rightarrow (2,2): & \quad Q_1((1,2), \text{\lr{right}}) = -1 + (0.1) [0 + 0.9\max\{0,0,0,0\} - (-1)] = -0.9 \\
    (2,2) \rightarrow (2,3): & \quad Q_1((2,2), \text{\lr{right}}) = 0 + (0.1) [+10 + 0.9\max\{0,0,0,0\} - 0] = +1
\end{align*}
در گام اول از مسیر دوم که مقدار $Q_1((1,1), \text{\lr{right}})$ را برای بار دوم محاسبه می‌کنیم مقدار متفاوتی نسبت به الگوریتم \lr{SARSA} تولید می‌شود ($-0.09$ در \lr{SARSA} و $0$ در \lr{Q-Learning}). این اتفاق به این علت رخ می‌دهد که در \lr{Q-Learning} هنگام محاسبه‌ی ارزش جدید، بین تمام اکشن‌ها بیشینه می‌گیریم. یعنی در \lr{SARSA} مقدار جدید بر اساس اکشن انجام شده محاسبه می‌شود، و در \lr{Q-Learning} بر اساس بهترین اکشن موجود.

\item با توجه به اینکه هنگامی که محرک وارد یک خانه می‌شود پاداش مربوط به آن را دریافت می‌کند، اگر پاداش هر گام را مشخص کنیم مسیرها به این صورت خواهند بود:
\begin{align*}
     & (1,1) \xrightarrow{0} (1,2) \xrightarrow{-10} (1,3)                       \\
     & (1,1) \xrightarrow{0} (1,2) \xrightarrow{0} (2,2) \xrightarrow{+10} (2,3)
\end{align*}
بنابراین ارزش این خانه‌ها به این صورت محاسبه می‌شود:
\begin{align*}
     & V^{\pi}(1,1) = \frac{(0 + 0.9(-10)) + (0 + 0.9(0) + 0.9^2(10))}{2} = \frac{-9 + 8.1}{2} = -0.45 \\
     & V^{\pi}(1,2) = \frac{(-10) + (0 + 0.9(+10))}{2} = \frac{-10 + 9}{2} = -0.5                      \\
     & V^{\pi}(2,2) = +10
\end{align*}
هیچ تخمینی برای ارزش خانه‌هایی که هیچ‌گاه استیت آغازین یا میانی مسیری نبوده‌اند نداریم. (در اینجا خانه‌های (1,3)، (2,1) و (2,3)) به طور کلی در روش \lr{Direct Evaluation}، هیچ اطلاعاتی بین همسایه‌ها منتقل نمی‌شود و ممکن است برای بعضی از خانه‌ها تخمینی نداشته باشیم. همچنین چون اطلاعات را ترکیب نمی‌کنیم، باید داده‌های بسیار زیادی داشته باشیم تا مقدار تخمینی از ارزش خانه‌ها به مقدار حقیقی نزدیک شود.

در مقابل در روش \lr{TD Learning}، برای به‌روز کردن ارزش خانه‌ها از خانه‌های همسایه نیز استفاده می‌کنیم. بنابراین این روش می‌تواند زودتر و با نمونه‌های کمتر، تخمین بهتری به دست آورد.

\end{ابجد}